Large language models can write plausible scientific ideas, but scientific ideation requires more than plausible novelty claims: ideas must be different from prior work and realistic enough to test. We examine whether novelty improves when it is grounded through concrete prompt operations rather than requested abstractly. Using 12 matched \aibidea contexts, we generate 84 ideas with \gptmini under seven conditions: an \ungrounded baseline plus deduction, induction, proposal writing, outcome imagination, a small lexical audit, and explicit literature comparison. We evaluate each idea with an independent \geminiflash judge on novelty, feasibility, specificity, grounding, risk awareness, and overall quality, and we add automatic overlap and diversity metrics. The result is mixed-to-negative. \litcompare has the highest mean overall score, 4.00 versus 3.92 for \ungrounded, but no judged quality metric shows a significant condition effect; for overall score, the Friedman test gives $p=0.157$. The \ungrounded baseline has the best \nfprod, 14.67. Grounded prompts do reliably change surface relation to prior work: prior-work overlap increases for every grounded condition, with a strong omnibus effect ($p=1.06\times 10^{-9}$). These results suggest that prompt-level grounding alone makes ideas sound more literature-connected without reliably improving idea quality.
